\section*{Capítulo \ref{ch:g2:animals-grow-up} --- Como os animais crescem}

\begin{solution}{exo:g2:animals-grow-up:1}
De um \omterm{def:g2:animals-grow-up:born}{ovo} posto pela mãe (um pintinho, um sapo, um caracol), ou de
barriga, saindo do corpo da mãe e bebendo o \omterm{def:g2:animals-grow-up:born}{leite} dela (um gatinho,
um bezerro, um bebê humano).
\end{solution}

\begin{solution}{exo:g2:animals-grow-up:2}
(a) De um \omterm{def:g2:animals-grow-up:born}{ovo}. (b) De barriga. (c) De um \omterm{def:g2:animals-grow-up:born}{ovo}. (d) De barriga. (e) De
um \omterm{def:g2:animals-grow-up:born}{ovo}.
\end{solution}

\begin{solution}{exo:g2:animals-grow-up:3}
O \omterm{def:g2:animals-grow-up:born}{leite} da mãe. Todos os \omterm{def:g1:animals-around-us:animal}{animais} que \omterm{def:g2:animals-grow-up:born}{nascem de barriga} neste
capítulo --- filhotes de cachorro, bezerros, coelhinhos, bebês
humanos --- também começam no \omterm{def:g2:animals-grow-up:born}{leite}.
\end{solution}

\begin{solution}{exo:g2:animals-grow-up:4}
A \omterm{def:g2:animals-grow-up:metamorphosis}{metamorfose} é a grande mudança de forma do corpo por que passam
alguns filhotes para virar adultos. O girino virando sapo e a
\omterm{ex:g2:animals-grow-up:butterfly}{lagarta} virando borboleta.
\end{solution}

\begin{solution}{exo:g2:animals-grow-up:5}
\omterm{def:g2:animals-grow-up:born}{Ovos}, girino, girino com \omterm{def:g1:my-body:parts}{pernas}, sapo.
\end{solution}

\begin{solution}{exo:g2:animals-grow-up:6}
\omterm{def:g2:animals-grow-up:born}{Ovo}, \omterm{ex:g2:animals-grow-up:butterfly}{lagarta}, \omterm{ex:g2:animals-grow-up:butterfly}{crisálida}, borboleta. É a \omterm{ex:g2:animals-grow-up:butterfly}{lagarta} que come \omterm{def:g1:plants-around-us:parts}{folhas} e
faz quase todo o crescimento.
\end{solution}

\begin{solution}{exo:g2:animals-grow-up:7}
O potrinho --- um cavalo pequeno desde o primeiro dia. Para ele,
crescer quer dizer ficar maior e mais forte mantendo a mesma forma.
\end{solution}

\begin{solution}{exo:g2:animals-grow-up:8}
Para mantê-los quentinhos. Dentro de cada \omterm{def:g2:animals-grow-up:born}{ovo} quentinho um pintinho
está sendo montado, alimentado pela gema, até estar pronto para
quebrar a casca e sair.
\end{solution}

\begin{solution}{exo:g2:animals-grow-up:9}
É um girino no meio da \omterm{def:g2:animals-grow-up:metamorphosis}{metamorfose}: as \omterm{def:g1:my-body:parts}{pernas} de trás já cresceram e
as da frente ainda não apareceram. Antes, ele era um girino sem
\omterm{def:g1:my-body:parts}{pernas} (e, antes disso, um \omterm{def:g2:animals-grow-up:born}{ovo}); depois, as \omterm{def:g1:my-body:parts}{pernas} da frente vão
aparecer, o rabo vai encolher e ele vai virar um sapinho.
\end{solution}

\begin{solution}{exo:g2:animals-grow-up:10}
Ninguém se mudou para lá --- a \omterm{ex:g2:animals-grow-up:butterfly}{lagarta} e a borboleta são o mesmo
\omterm{def:g1:animals-around-us:animal}{animal}. Dentro da \omterm{ex:g2:animals-grow-up:butterfly}{crisálida}, o corpo da \omterm{ex:g2:animals-grow-up:butterfly}{lagarta} foi reconstruído
como o de uma borboleta: essa reconstrução é o que a \omterm{def:g2:animals-grow-up:metamorphosis}{metamorfose}
quer dizer.
\end{solution}
